\documentclass[12pt]{article}
\usepackage[utf8]{inputenc}
\usepackage[T1]{fontenc}
\usepackage{hyperref}
\title{Cayuse Language Resources for K-12 Curriculum Development}
\author{Compiled by Claude Sonnet 5.5 (claude-sonnet-5-5) for Andr\'as Kornai}
\date{}
\begin{document}
\maketitle

\section{Introduction}
Cayuse (Weyíiletpuu; ISO 639-3 xcy; Glottocode cayu1241) was the language of the Cayuse people of Washington and Oregon. It is usually treated as an isolate. The Confederated Tribes of the Umatilla Indian Reservation (CTUIR) language pages describe it differently, as a dialect of Nez Perce that has not been used since the 1940s. CTUIR-linked sources variously call it dormant or extinct. Revitalization work at CTUIR focuses on Umatilla and Walla Walla, and I found no Cayuse-specific school programme.

\section{Primary Educational Resources}
None specific to Cayuse were found. The CTUIR language department page \url{https://ctuir.org/departments/education/language/} covers the tribal language programme. The immersion work I saw described was for Umatilla, not Cayuse.

\section{Linguistic and Reference Materials}
The CTUIR online dictionary has an entry for the Cayuse people, giving the form /weyíletpuu/: \url{https://dictionary.ctuir.org/?p=46373}. I did not find a Cayuse grammar or dictionary.

\section{Religious and Community Resources}
CTUIR is the community organization. See \url{https://ctuir.org/about/history-culture/} and \url{https://ctuir.org/about/language/}. A tribal newspaper article, \url{https://cuj.ctuir.org/?p=4308}, discusses language loss. It reports that a Umatilla teacher has struggled to work from a small document of Cayuse words. Another page, \url{https://www.legendsofamerica.com/cayuse-tribe/}, gives general history.

\section{Digital and Audio Resources}
No Cayuse audio was found. The CTUIR dictionary site above is for Umatilla and related languages.

\section{Curriculum Development Context}
Very little Cayuse material exists. A realistic K-12 programme would combine Cayuse history and heritage with the Umatilla language materials of CTUIR. Any Cayuse-language content would need to be developed with the CTUIR Language Program, whose authority on the language's status should be respected.
\end{document}